\documentclass[11pt]{article}
\usepackage[margin=1in]{geometry}
\usepackage{amsmath,amssymb,amsthm,hyperref}
\newtheorem{theorem}{Theorem}
\newtheorem{lemma}{Lemma}
\title{The Guo--Liu bound for weighted Hermite polynomials}
\author{Principia Math In Progress}
\date{17 September 2026}
\begin{document}
\maketitle
\noindent MathDB \href{https://mathdb.com/p/346243}{346243};
Guo--Liu, Conjecture 1, equation (29) in the published paper
\cite{GL}, and equation (30) in arXiv:2004.10441v1.

\section*{Task 1. Prove the original bound for every degree}
Use the physicists' normalization
\[
 H_n(x)=(-1)^ne^{x^2}\frac{d^n}{dx^n}e^{-x^2}.
\]
Write
\[
 N_n=2n+1,\quad C_n=\frac{2^n\Gamma((n+1)/2)}{\sqrt\pi},\quad
 g_n(x)=(1-x^2/N_n)^{-1/2},\quad
 x_n=\sqrt{N_n}\left(1-\frac{\pi}{2N_n}\right).
\]
In particular $0<x_n<\sqrt{N_n}$ for every $n\ge1$.

\begin{theorem}\label{thm:main}
For every integer $n\ge1$ and every real $x>0$,
\begin{equation}\label{eq:target}
 |H_n(x)e^{-x^2}|
 \le C_n g_n(\min\{x,x_n\})e^{-x^2/2}.
\end{equation}
\end{theorem}

Put $y_n(x)=H_n(x)e^{-x^2/2}$. The two standard identities needed below
follow directly from the Hermite differential equation and
$H_{n+1}'=2(n+1)H_n$:
\begin{equation}\label{eq:ode}
 y_n''+(N_n-x^2)y_n=0,\qquad
 y_{n+1}'+xy_{n+1}=2(n+1)y_n.
\end{equation}
We also use the classical global Hermite inequality
\begin{equation}\label{eq:global}
 |y_n(x)|\le\sqrt{2^n n!}\qquad(x\in\mathbb R),
\end{equation}
in precisely the normalization recorded in DLMF 18.14.9 \cite{DLMF}.
This known inequality is an input, not a new assertion of this note.

\subsection*{The local bound, including the odd degrees}
\begin{lemma}\label{lem:local}
For every integer $n\ge1$ and $0\le x<\sqrt{N_n}$,
\[
 |y_n(x)|\le C_n(1-x^2/N_n)^{-1/2}.
\]
\end{lemma}
\begin{proof}
Define the oscillator energy
\[
 E_n(x)=y_n'(x)^2+(N_n-x^2)y_n(x)^2.
\]
Equation \eqref{eq:ode} gives $E_n'(x)=-2xy_n(x)^2\le0$ for $x\ge0$.
For even $n$, $y_n'(0)=0$ and
$|y_n(0)|=n!/(n/2)!=C_n$, so
\[
 (N_n-x^2)y_n(x)^2\le E_n(x)\le N_n C_n^2.
\]
This proves the assertion for even degrees.

For odd $n$, apply the even-degree energy estimate to $v=y_{n+1}$.
Set $M=N_n+2$ and $q=M-x^2>0$. Weighted Cauchy--Schwarz and
\eqref{eq:ode} yield
\begin{align*}
 4(n+1)^2y_n(x)^2
   &=(v'(x)+xv(x))^2\\
   &\le (v'(x)^2+qv(x)^2)(1+x^2/q)\\
   &\le M C_{n+1}^2\frac{M}{M-x^2}.
\end{align*}
Thus, in fact, throughout $0\le x<\sqrt M$,
\begin{equation}\label{eq:oddstrong}
 |y_n(x)|\le D_n(1-x^2/M)^{-1/2},\qquad
 D_n=\frac{C_{n+1}\sqrt M}{2(n+1)}.
\end{equation}
We claim $D_n<C_n$. Write $n=2m-1$, $m\ge1$. The gamma duplication
formula gives
\[
 \left(\frac{D_n}{C_n}\right)^2
 =Q_m:=\left(m+\frac14\right)
       \left(\frac{\Gamma(m+1/2)}{\Gamma(m+1)}\right)^2.
\]
Its consecutive ratio satisfies the exact identity
\[
 \frac{Q_{m+1}}{Q_m}
 =1+\frac1{4(m+1)^2(4m+1)}>1.
\]
Stirling's formula gives $Q_m\longrightarrow1$. A strictly increasing
sequence with limit $1$ has every term less than $1$, proving the claim.
Finally, for $x^2<N_n<M$,
\[
 (1-x^2/M)^{-1/2}\le(1-x^2/N_n)^{-1/2}.
\]
Together with \eqref{eq:oddstrong}, this completes the odd case.
\end{proof}

\subsection*{The capped region for degrees at least three}
Let $A_n=C_ng_n(x_n)$. Its exact value satisfies
\[
 A_n^2=C_n^2\frac{N_n^2}{\pi(N_n-\pi/4)}.
\]
We show $A_n^2>2^nn!$ for all $n\ge3$. Define
\[
 r_n=\frac{C_n^2}{2^nn!}
 =\frac{\Gamma((n+1)/2)}{\sqrt\pi\,\Gamma(n/2+1)},\qquad
 U_n=\frac{A_n^2}{2^nn!}
 =\frac{r_nN_n^2}{\pi(N_n-\pi/4)}.
\]
Since $r_{n+2}/r_n=(n+1)/(n+2)$, putting $N=N_n$ and $a=\pi/4<1$
gives
\[
 \frac{U_{n+2}}{U_n}
 =\frac{(N+1)(N+4)^2(N-a)}{(N+3)N^2(N+4-a)}>1.
\]
For completeness, the numerator minus the denominator is
\begin{align*}
 &2\{N^3+(6-3a)N^2+(8-12a)N-8a\}\\
 &\hspace{1cm}>2(N^3+3N^2-4N-8)>0\qquad(N\ge3).
\end{align*}
The last inequality follows from $N^3-4N>0$ and $3N^2-8>0$.
The elementary bound $\pi<22/7$, and the monotonicity of the relevant
denominators on this interval, give the two base cases
\[
 U_3=\frac{784}{3\pi^2(28-\pi)}
     >\frac{33614}{31581}>1,\qquad
 U_4=\frac{243}{2\pi(36-\pi)}
     >\frac{11907}{10120}>1.
\]
Induction in each parity proves the claim. Therefore \eqref{eq:global}
implies $|y_n(x)|<A_n$ for every real $x$ when $n\ge3$.

\subsection*{The two small degrees}
Direct differentiation gives
\[
 \sup_{x\ge0}|y_1(x)|^2=4/e,\qquad
 \sup_{x\ge0}|y_2(x)|^2=\max\{4,64e^{-5/2}\}.
\]
Here $H_1(x)=2x$ and $H_2(x)=4x^2-2$; the respective nonzero
critical points are $x=1$ and $x=\sqrt{5/2}$.
Again using $\pi<22/7$,
\[
 A_1^2=\frac{144}{\pi^2(12-\pi)}>
 \frac{6174}{3751}>\frac85>\frac4e,
\]
because $e>5/2$. Similarly,
\[
 A_2^2=\frac{400}{\pi(20-\pi)}>
 \frac{4900}{649}>\frac{15}{2}>\max\{4,64e^{-5/2}\}.
\]
For the last step it suffices that
$e^{5/2}>1+5/2+25/8+125/48=443/48>9$, since $64/9<15/2$.
All comparisons are exact, and give $|y_n(x)|<A_n$ also for $n=1,2$.

\paragraph{Completion of the proof.}
For $0<x\le x_n$, apply Lemma \ref{lem:local}. For $x\ge x_n$, apply
the global cap just established for every degree. Multiplication by
$e^{-x^2/2}$ proves \eqref{eq:target} in both regions.\hfill$\square$

\paragraph{Scope.}
The proof covers every integer degree, every positive real argument,
and the exact cutoff prescribed in the source. It also proves the
stronger odd-degree local estimate \eqref{eq:oddstrong}. Exact symbolic
checks accompany the paper; finite numerical samples are not used to
justify the universal inequality. The proof was reviewed by the
authoring agent; independent peer review is not claimed.

\section*{Task 2. Optimize the common cutoff constant}
Replace the source's $\pi/2$ by a constant $c\in(0,3)$ in
$x_n(c)=\sqrt{N_n}(1-c/N_n)$. Determine the largest $c$ for which the
capped bound holds for all $n\ge1$. The degree-one maximum gives an
explicit necessary restriction; determine whether that restriction is
also sufficient for every higher degree.

\section*{Task 3. Optimize the local envelope}
For each degree, determine the least exponent $\theta_n$ such that
$|y_n(x)|\le C_n(1-x^2/N_n)^{-\theta_n}$ throughout
$0\le x<\sqrt{N_n}$. Study its large-degree limit, its parity dependence,
and the locations at which the optimal bound is attained. Compare the
result with the stronger odd-degree estimate in this proof.

\section*{Task 4. Improve the summation remainder quantitatively}
Insert the minimum of the local energy bound and the classical global
bound into the one- and two-dimensional remainder integrals of
Guo--Liu, instead of using the constant $A_n$ throughout. Derive explicit
computable remainder bounds and optimize the choices of truncation
rank and Fourier cutoff for a prescribed tolerance.

\section*{Task 5. Extend the parity transfer to generalized Hermite functions}
For the generalized Hermite family associated with
$|x|^{2\mu}e^{-x^2}$, $\mu>-1/2$, determine which adjacent-degree
differential identities yield an analogous local bound with explicit
gamma constants. Establish the admissible parameter range and a
uniform transition to a finite global cap, accounting for the singular
term at the origin in the transformed differential equation.

\begin{thebibliography}{9}
\bibitem{GL}
J. Guo and Y. Liu, \emph{A modified Euler--Maclaurin formula in 1D and
2D with applications in statistical physics}, Commun. Theor. Phys.
\textbf{73} (2021), 075002, Conjecture 1, equation (29), p. 8.
\href{https://doi.org/10.1088/1572-9494/abf1dd}{doi:10.1088/1572-9494/abf1dd}.
Preprint: \href{https://arxiv.org/abs/2004.10441}{arXiv:2004.10441v1}.
\bibitem{DLMF}
NIST Digital Library of Mathematical Functions,
\href{https://dlmf.nist.gov/18.14.E9}{equation 18.14.9},
global bound for Hermite polynomials; accessed 17 September 2026.
\end{thebibliography}
\end{document}
